\section{Introduction}

Introduce the research problem and explain why it matters.
This opening paragraph should define the application setting, the current limitations of prior work, and the practical or scientific motivation for the paper.
Use concrete language, but avoid implementation details that belong in later sections.

Summarize the central gap addressed by the paper.
For example, the gap may involve quality, efficiency, robustness, scalability, interpretability, data availability, or deployment constraints.
Frame the gap in a way that prepares the reader for the proposed method.

Briefly describe the proposed approach.
This paragraph should give the reader a high-level mental model of the system, algorithm, dataset, or analysis before the technical sections begin.
Refer to the main overview figure when appropriate, such as Figure~\ref{fig:teaser}.

The main contributions can be listed as follows:
\begin{itemize}
    \item \textbf{Contribution One.} Describe the first technical or empirical contribution in one or two sentences.
    \item \textbf{Contribution Two.} Describe the second contribution, emphasizing what is new or useful.
    \item \textbf{Contribution Three.} Describe the validation, dataset, benchmark, application, or analysis contribution.
\end{itemize}

Close the introduction with a short summary of the evidence supporting the paper's claims.
Mention the evaluation setting, important metrics, and broad conclusions without including final numeric results until they are ready for release.

\newpage
